
% !TEX root = main.tex
\appendices
\section{Ethics}\label{sec:ethics}
%\jz{Potentially move to appendix for space. CFP: All papers must include a statement about ethical issues in the appendix. }
We follow best practices for measurements~\cite{ZMapInternetScanner2013}, and we have obtained ethical approval from our institution.
No sensitive information is collected.
We received one complaint from our \gls{cert} due to high scanning traffic.
We limited the scanning rate by distributing the pipeline over a 24-hour period and used a blocklist from prior measurements.

\section{IPv6 Anycast}\label{sec:ipv6-anycast}

\begin{table}[H]
    \centering
    \caption{Top 10 ASes by number of anycast IPv6 /48-prefixes (n=\num{13007})}
    \begin{tabular}{l r@{~}l}
        \toprule
        Org (ASN) & \multicolumn{2}{c}{/48-Prefixes} \\
        \midrule
        Cloudflare Spectrum (209242) & \num{6421} & (49\%) \\
        Cato Networks (214040) & \num{1795} & (13\%) \\
        Fastly (54113) & \num{760} & (6\%) \\
        Cloudflare (13335) & \num{620} & (5\%) \\
        Incapsula (19551) & \num{587} & (5\%) \\
        \bottomrule
    \end{tabular}
    \label{tab:top_ases_v6}
\end{table}

To understand the ASes that deploy the most IPv6 anycast services, we present~\cref{tab:top_ases_v6}.
Unlike IPv4, we do not know the number of active IPs for IPv6 due to the large number of addresses within /48-prefixes.
However, based on the number of /48s covered by \laces, we find that Cloudflare is, by far, the largest operator of IPv6 anycast.
We note that \laces has limited IPv6 visibility, as their hitlists do not cover the full IPv6 address space.

\section{Anycast upstreams}\label{sec:anycast-upstreams}

\begin{table}[H]
    \centering
    \caption{Top 5 upstreams visible for routed anycast prefixes (n=7,442)}
    \begin{tabular}{crrr}
        \toprule
        Name (ASN) & Total Prefixes & IPv4 Prefixes & IPv6 Prefixes \\
        \midrule
        \textbf{Arelion} (1299) & \num{2977} (40\%) & \num{1751} (37\%) & \num{1226} (45\%) \\
        \textbf{Cogent} (174) & \num{2577} (35\%) & \num{1557} (33\%) & \num{1020} (37\%) \\
        \textbf{NTT} (2914) & \num{2153} (28\%) & \num{1171} (25\%) & \num{982} (36\%) \\
        \textbf{Lumen} (3356) & \num{2049} (28\%) & \num{1485} (32\%) & \num{564} (21\%) \\
        \textbf{Telxius} (12956) & \num{1557} (21\%) & \num{1076} (23\%) & \num{481} (17\%) \\
        %\textbf{TATA} (6453) & 1,531 (20.6\%) & 1,150 (24.5\%) & 381 (13.8\%) \\
        %\textbf{Telstra} (4637) & 1,323 (17.8\%) & 624 (13.3\%) & 699 (25.4\%) \\
        %\textit{NetActuate} (36236) & 927 (12.5\%) & 554 (11.8\%) & 373 (13.5\%) \\
        %\textbf{GTT} (3257) & 750 (10.1\%) & 529 (11.3\%) & 221 (8.0\%) \\
        %\textbf{Vocus} (4826) & 708 (9.5\%) & 402 (8.6\%) & 306 (11.1\%) \\
        %\textbf{PCCW} (3491) & 555 (7.5\%) & 378 (8.1\%) & 177 (6.4\%) \\
        %\textit{Vultr} (20473) & 482 (6.5\%) & 249 (5.3\%) & 233 (8.5\%) \\
        %\textit{Cloudflare} (13335) & 415 (5.6\%) & 368 (7.9\%) & 47 (1.7\%) \\
        %\textbf{RETN} (9002) & 372 (5.0\%) & 160 (3.4\%) & 212 (7.7\%) \\
        %\textbf{AS6762} & Sparkle & 302 (4.1\%) & 151 (3.2\%) & 151 (5.5\%) \\
        %\textbf{AS8455} & atom86 & 284 (3.8\%) & 142 (3.0\%) & 142 (5.2\%) \\
        %\textit{AS20940} & Akamai & 219 (2.9\%) & 135 (2.9\%) & 84 (3.0\%) \\
        %\textbf{AS6461} & Zayo & 200 (2.7\%) & 185 (3.9\%) & 15 (0.5\%) \\
        %\textit{AS15169} & Google & 192 (2.6\%) & 178 (3.8\%) & 14 (0.5\%) \\
        %\textbf{AS6939} & HE & 186 (2.5\%) & 49 (1.0\%) & 137 (5.0\%) \\
        \bottomrule
    \end{tabular}
    \label{tab:anycast_upstreams}
\end{table}


Table~\ref{tab:anycast_upstreams}
shows the most frequent upstream ASes for routed prefixes covering the /24, and /48 anycast prefixes of \laces.
%
Unsurprisingly, these are all tier-1 networks (except Cogent for IPv6). % https://en.wikipedia.org/wiki/Tier_1_network
%
Comparing IPv4 and IPv6, we find they provide roughly equal connectivity.

\section{Multi-Regional Providers}\label{sec:regional-anycast}

\begin{table}[H]
    \centering
    \caption{ASes with multiple regional anycast deployments inside different continents (n=\num{23})}
    \begin{tabular}{lr}
        \toprule
        Org (ASN) & Regions (\# of /24s) \\
        \midrule
        AWS (16509) & AS (11), NA (7), EU (1), OC (1), SA (1) \\
        BunnyCDN (200325) & NA (2), AS (1), OC (1), SA (1) \\
        Akamai (33905) & AS (10), NA (10), OC (10), EU (1) \\
        Fortanix (23013) & AS (1), EU (1), NA (1) \\
        GCORE (199524) & EU (2), NA (2), AS (1) \\
        Vultr (20473) & NA (3), AS (1), OC (1) \\
        Fortinet (40934) & NA (5), AS (2), EU (2) \\
        Bytedance (396986) & EU (3), AS (2), NA (1) \\
        Cato networks (13150) & NA (5), EU (3), AS (1) \\
        Level 3 (3356) & NA (8), EU (1) \\
        Fastly (54113) & NA (180), EU (1) \\
        Radware (48851 \& 25773) & EU (18), NA (14), OC (1) \\
        Radware (48851) & EU (18), OC (1) \\
        F5 (35280) & NA (2), EU (1) \\
        RIPE-NCC-RIS-AS (12654) & AS (1), NA (1) \\
        Voxility (3223) & EU (1), NA (1) \\
        Cache Networks (30081) & AS (1), EU (1) \\
        Akamai (21342) & NA (4), EU (1) \\
        FLASHSTART (204903) & EU (1), NA (1) \\
        Radware (198949) & EU (2), NA (2) \\
        ACE (139341) & NA (3), AS (2) \\
        Huawei Clouds (136907) & AS (2), EU (2) \\
        Netactuate (63911) & NA (16), EU (1) \\
        \bottomrule
    \end{tabular}
    \label{tab:anycast_multi_continent}
\end{table}

Table~\ref{tab:anycast_multi_continent} lists all ASes for which \laces detects regional (\ie, confined within a single continent) anycast prefixes in multiple continents.
We note that the iGreedy~\cite{igreedy} geolocation method used by \laces struggles with detecting regional anycast and may miss certain prefixes.

\section{DNS Software}\label{app:dns-software}
We request \texttt{TXT CHAOS} records for commonly used records such as \texttt{id.server} and \texttt{version.bind} to investigate the DNS software running on anycast prefixes.
%
For example, \textit{ns.es.hk2.bind} strongly indicates that this particular DNS server in Hong Kong uses BIND\@.
%
We perform this measurement using Ark~\cite{ark} that lets us request these records from 350+ different \glspl{vp}.

In total, \num{658} ASes run an anycasted DNS service
and \num{190} ASes (29\%) reveal the software used in the requested TXT records.
%
As observed by Cicalese et al.~\cite{CicaleseEtAl_CharacterizingIPv4Anycast_2015}, we find BIND is the most commonly used DNS software, seen in \num{92} ASes, followed by PowerDNS with \num{86} ASes.

Interestingly, we find \num{25} ASes that deploy multiple DNS software using the same anycast IP address, where, depending on the \gls{pop} reached, a different DNS software is observed.
%
Examples include: K- and B-root, AS112, AWS, Google, Verizon, CIRA (.ca), and RIR (LACNIC, RIPE, APNIC) DNS infrastructure such as rDNS\@.
